\chapter{Macroéconomie, connaissance des marchés et pitches}\label{ch:iv2:macro-market-knowledge-and-pitches}

\begin{ivtrains}
Comment les taux, l'inflation et la croissance font bouger les actifs, ce que fait un événement aux prix une fois qu'on se demande ce qui était déjà intégré, comment suivre les niveaux, et comment préparer un pitch d'action ou une idée de trade. Tous les niveaux donnés ici sont inventés pour le calcul. Avant tout entretien, il faut les remplacer par ceux du jour, tirés d'une page dont on peut citer la source.
\end{ivtrains}

\begin{question}[$\star$ \iqroles{Analyste, Trader, QT} \iqfirm{toutes} \iqlevel{junior}]\label{iq:iv2:macro-market-knowledge-and-pitches:1}
Les rendements montent de 1 point lors d'une surprise. Que devient une obligation de duration 7~? Et une action valorisée $D/(r - g)$ avec $r = 8~\%$ et $g = 4~\%$~? Et si les anticipations de croissance montent aussi d'un demi-point~?
\end{question}

\begin{question}[$\star$ \iqroles{Analyste, Trader, QT} \iqfirm{toutes} \iqlevel{junior}]\label{iq:iv2:macro-market-knowledge-and-pitches:2}
Une obligation d'État nominale à cinq ans affiche un rendement de 4.2~\%, et l'obligation indexée sur l'inflation rapporte 1.7~\%. Quel est le point mort d'inflation~? Que signifie-t-il, et qu'est-ce qu'il contient d'autre~?
\end{question}

\begin{question}[$\star$ \iqroles{Analyste, Trader, QR} \iqfirm{toutes} \iqlevel{junior}]\label{iq:iv2:macro-market-knowledge-and-pitches:3}
Le chiffre d'inflation publié ressort au-dessus du consensus. Qu'est-ce qui bouge, dans quel sens et pourquoi~? Passez en revue les taux courts, les taux longs, la devise, les actions et l'or.
\end{question}

\begin{question}[$\star$ \iqroles{Analyste, Trader, QT} \iqfirm{toutes} \iqlevel{junior}]\label{iq:iv2:macro-market-knowledge-and-pitches:4}
Le taux directeur est de 4.33~\%, une réunion approche, et les contrats à terme sur la période qui suit cette réunion affichent 4.20~\%. Quelle probabilité de baisse de 25 points de base est intégrée dans les cours~? La banque centrale maintient ses taux. Que deviennent les taux courts et les actions, et pourquoi la phrase «~la banque centrale a maintenu~» ne vous dit-elle rien à elle seule~?
\end{question}

\begin{question}[$\star$ \iqroles{Analyste, Trader} \iqfirm{toutes} \iqlevel{junior}]\label{iq:iv2:macro-market-knowledge-and-pitches:5}
Une société présente ses comptes en dollars, et 80~\% de ses ventes se font en Europe. L'euro recule de 10~\% face au dollar, et les ventes en monnaie locale ne changent pas. Que devient le chiffre d'affaires publié~? Que demanderiez-vous ensuite~?
\end{question}

\begin{question}[$\star$ \iqroles{Analyste, Trader, QT, QR} \iqfirm{toutes} \iqlevel{junior, confirmé}]\label{iq:iv2:macro-market-knowledge-and-pitches:6}
Votre recruteur vous demande les niveaux actuels du rendement à dix ans, de l'indice actions principal, du pétrole et des principales paires de devises. Comment vous préparez-vous à cela, et que répondez-vous si vous ne le savez pas~?
\end{question}

\begin{question}[$\star$ \iqroles{Analyste, Trader} \iqfirm{toutes} \iqlevel{junior}]\label{iq:iv2:macro-market-knowledge-and-pitches:7}
Une compagnie aérienne consomme 1 milliard de gallons de carburant par an et réalise un résultat d'exploitation de 1~milliard de dollars. Le brut monte de 10~dollars le baril, et un baril contient 42 gallons. Quel est l'ordre de grandeur de l'impact sur le résultat~? Que changerait le fait que la moitié du carburant soit couverte~?
\end{question}

\begin{question}[$\star\star$ \iqroles{Analyste, Trader, QR} \iqfirm{toutes} \iqlevel{junior, confirmé}]\label{iq:iv2:macro-market-knowledge-and-pitches:8}
Le rendement à deux ans est de 4.5~\%, et celui à dix ans de 4.0~\%, avec capitalisation annuelle. Quel taux à huit ans, démarrant dans deux ans, cela implique-t-il~? Une courbe inversée annonce-t-elle une récession~? Que dit-elle exactement~?
\end{question}

\begin{question}[$\star\star$ \iqroles{Analyste, Trader} \iqfirm{toutes} \iqlevel{junior, confirmé}]\label{iq:iv2:macro-market-knowledge-and-pitches:9}
Les estimations de bénéfices montent de 5~\% grâce à une économie plus forte, tandis que le taux d'actualisation monte de 50 points de base. L'action a une duration de 20. Quel est l'effet net~? Quand les actions et les rendements montent-ils ensemble~?
\end{question}

\begin{question}[$\star\star$ \iqroles{Analyste, Trader, QR} \iqfirm{banque, hedge fund} \iqlevel{junior, confirmé}]\label{iq:iv2:macro-market-knowledge-and-pitches:10}
Une règle de politique monétaire s'écrit $i = r^* + \pi + 0.5(\pi - \pi^*) + 0.5y$, avec $r^* = 2~\%$, $\pi^* = 2~\%$, une inflation de 3~\% et un écart de production de 1~\%. Quel taux en résulte-t-il~? Que se passe-t-il si l'inflation monte d'un point~? Que ne peut pas vous dire une règle de ce type~?
\end{question}

\begin{question}[$\star\star$ \iqroles{QR, Risque, Analyste} \iqfirm{toutes} \iqlevel{junior, confirmé}]\label{iq:iv2:macro-market-knowledge-and-pitches:11}
Un portefeuille 60/40 détient des actions avec une volatilité de 16~\% et des obligations avec une volatilité de 7~\%. Quelle est sa volatilité avec une corrélation actions-obligations de $-0.3$, puis de $+0.3$~? Qu'est-ce qui fait basculer le signe de la corrélation~?
\end{question}

\begin{question}[$\star\star$ \iqroles{Trader, QT, Analyste} \iqfirm{banque, hedge fund} \iqlevel{junior, confirmé}]\label{iq:iv2:macro-market-knowledge-and-pitches:12}
Les créations d'emplois hors agriculture dépassent le consensus de 150\,000. D'après les données passées, les rendements à deux ans bougent de 1 point de base pour 30\,000 de surprise. Quel est le mouvement attendu~? Combien perdez-vous avec un DV01 de 10\,000~dollars en position longue~? Pourquoi les rendements pourraient-ils rester stables, voire baisser~?
\end{question}

\begin{question}[$\star\star$ \iqroles{Trader, Analyste, QR} \iqfirm{banque, hedge fund} \iqlevel{junior, confirmé}]\label{iq:iv2:macro-market-knowledge-and-pitches:13}
Une banque centrale baisse ses taux de 50 points de base alors que 25 étaient intégrés dans les cours. D'après les données passées, une surprise se transmet à hauteur de 60~\% aux rendements à deux ans et de 20~\% à ceux à dix ans. Que se passe-t-il pour chacun, et pour un steepener de 10\,000~dollars de DV01, acheteur sur deux ans et vendeur sur dix ans~? Que pourrait faire le reste du marché~?
\end{question}

\begin{question}[$\star\star$ \iqroles{Trader, QT, Analyste} \iqfirm{banque, hedge fund} \iqlevel{junior, confirmé}]\label{iq:iv2:macro-market-knowledge-and-pitches:14}
Vous empruntez des yens à 0.5~\% et investissez en dollars à 4.5~\%. De combien le dollar peut-il baisser face au yen en un an avant de perdre de l'argent~? Avec une volatilité de 10~\%, quel est le ratio de Sharpe~? Que se passe-t-il lors d'un débouclage~?
\end{question}

\begin{question}[$\star\star$ \iqroles{Analyste, Trader, QR} \iqfirm{banque, hedge fund} \iqlevel{junior, confirmé}]\label{iq:iv2:macro-market-knowledge-and-pitches:15}
Que fait une banque centrale lorsqu'elle achète ou vend des obligations à grande échelle, et comment cet effet se répercute-t-il sur le rendement~? Pourquoi cet effet est-il difficile à mesurer~?
\end{question}

\begin{question}[$\star\star$ \iqroles{Analyste, Trader} \iqfirm{toutes} \iqlevel{junior}]\label{iq:iv2:macro-market-knowledge-and-pitches:16}
Une société bat les attentes sur ses bénéfices de 5~\%, mais abaisse ses prévisions. Le P/E de l'action passe de 20 à 18. Que devient le cours~? Qu'est-ce qui importait au marché~?
\end{question}

\begin{question}[$\star\star$ \iqroles{Analyste, Trader, QT} \iqfirm{hedge fund, banque} \iqlevel{junior, confirmé}]\label{iq:iv2:macro-market-knowledge-and-pitches:17}
Faites-moi le pitch d'une action. Quelle structure utilisez-vous, que doit-elle contenir, et comment construisez-vous la vôtre~?
\end{question}

\begin{question}[$\star\star\star$ \iqroles{Analyste, Trader, QR} \iqfirm{hedge fund} \iqlevel{junior, confirmé}]\label{iq:iv2:macro-market-knowledge-and-pitches:18}
Vous présentez une action à 50~dollars. Scénario haussier~: BPA de 3.60 valorisé à un P/E de 18, avec une probabilité de 50~\%. Scénario de base~: 3.30 à 16, avec 30~\%. Scénario baissier~: 2.70 à 12, avec 20~\%. Quelle est l'espérance de rendement~? Quel est le rapport gain/risque~? Si vous prévoyez de sortir en stop à 44~dollars et que l'idée peut vous faire perdre 1~\% du portefeuille, quelle taille de position prenez-vous~? Que peut-il mal tourner avec ce stop~?
\end{question}

\begin{question}[$\star\star\star$ \iqroles{Trader, QT, Analyste} \iqfirm{banque, hedge fund} \iqlevel{junior, confirmé}]\label{iq:iv2:macro-market-knowledge-and-pitches:19}
Présentez un steepener 2s10s neutre en DV01~: acheteur sur deux ans, vendeur sur dix ans, 10\,000~dollars de DV01 de chaque côté. Les durations modifiées sont de 1.9 et 8. Les rendements sont de 4.0~\% et 4.2~\%, et le repo est à 4.1~\%. Quels sont les notionnels~? Combien gagnez-vous sur 25 points de base de steepening~? Quel est le carry annuel~? Avec une volatilité quotidienne du spread de 6 points de base, quelle est la VaR à un jour à 99~\%~?
\end{question}

\begin{question}[$\star\star\star$ \iqroles{Analyste, QR, Risque} \iqfirm{toutes} \iqlevel{junior, confirmé}]\label{iq:iv2:macro-market-knowledge-and-pitches:20}
L'inflation est de 3~\% par an. Une obligation indexée sur l'inflation et une obligation nominale ont chacune une duration de 8. L'obligation nominale rapporte 4~\%, et l'indexée un rendement réel de 1.7~\%. Comparez les rendements à un an si les rendements nominaux montent de 1 point à rendements réels inchangés, et si les rendements réels montent de 1 point à points morts d'inflation inchangés. Que couvre une obligation indexée sur l'inflation, et que ne couvre-t-elle pas~?
\end{question}

\begin{question}[$\star\star\star$ \iqroles{Analyste, Trader, QR} \iqfirm{banque, hedge fund} \iqlevel{junior, confirmé}]\label{iq:iv2:macro-market-knowledge-and-pitches:21}
Le rendement nominal à dix ans monte de 50 points de base en une semaine~: les rendements réels montent de 60, et les points morts d'inflation baissent de 10. Qu'est-ce qui a bougé, et quelle est l'explication~? Comparez les mouvements des obligations nominales de duration 7.5 et des obligations indexées de duration 8. Préféreriez-vous détenir des actions de croissance ou de valeur, et pourquoi~?
\end{question}

\begin{question}[$\star\star\star$ \iqroles{Analyste, QR, Trader} \iqfirm{toutes} \iqlevel{junior, confirmé}]\label{iq:iv2:macro-market-knowledge-and-pitches:22}
Présentez la grille croissance-inflation~: croissance à la hausse ou à la baisse, inflation à la hausse ou à la baisse. Que tendent à faire les actions, les obligations et les matières premières dans chaque case, et quels actifs gagnent~? Pourquoi un tableau de ce type est-il dangereux~?
\end{question}

\begin{question}[$\star\star\star$ \iqroles{QR, Risque, Analyste} \iqfirm{toutes} \iqlevel{junior, confirmé}]\label{iq:iv2:macro-market-knowledge-and-pitches:23}
Un investisseur américain détient des actions en euros, qui rapportent 10~\% en euros. L'euro recule de 5~\%. Les taux en dollars sont de 4.5~\% et les taux en euros de 2.5~\%. Quel est le rendement en dollars sans couverture~? Quel est le rendement avec une couverture d'un an sur la valeur de départ~? Pourquoi un écart de taux positif représente-t-il un gain pour cette couverture~?
\end{question}

\begin{question}[$\star\star\star$ \iqroles{Trader, QR, Analyste} \iqfirm{banque, hedge fund} \iqlevel{junior, confirmé}]\label{iq:iv2:macro-market-knowledge-and-pitches:24}
Le taux directeur est de 4.33~\%. Une réunion a lieu au milieu du mois 1, et le contrat à terme du mois 1 affiche en moyenne 4.25~\%. La réunion suivante a lieu le premier jour du mois 2, et ce contrat affiche en moyenne 4.10~\%. Une variation est de 25 points de base. Quelles probabilités de baisse ces niveaux impliquent-ils à chaque réunion~? Que faut-il ajuster en pratique~?
\end{question}

\begin{question}[$\star\star\star\star$ \iqroles{Trader, Analyste, QR, Risque} \iqfirm{hedge fund, banque} \iqlevel{confirmé}]\label{iq:iv2:macro-market-knowledge-and-pitches:25}
Les données d'inflation tombent demain, avec un consensus à 0.3~\% en variation mensuelle. Votre anticipation~: 0.2~\% avec une probabilité de 25~\%, 0.3~\% à 50~\%, 0.4~\% à 25~\%. Vous associez ces cas à des rendements d'actions de $+1~\%$, $+0.3~\%$ et $-1.5~\%$, et à des variations du rendement à deux ans de $-8$, 0 et $+10$ points de base. Vous détenez 100~millions de dollars d'actions et un DV01 acheteur de 50\,000~dollars sur des obligations à deux ans. Calculez le P\&L dans chaque cas, sa moyenne et son écart-type. Un straddle sur indice à la monnaie pour la journée coûte 1~\%. Vaut-il la peine de l'acheter~? Que faites-vous ce soir, et que diriez-vous à votre manager~?
\end{question}
